% ----------------------------------------------------------------------
%  Pracovní úkoly
% ----------------------------------------------------------------------
\section{Pracovní úkoly}

\begin{enumerate}
\item S využitím krystalu LiF jako analyzátoru proveďte měření následujících rentgenových spekter:

\begin{enumerate}
    \item Rentgenka s Cu anodou.

    \begin{enumerate}
        \item proměřte krátkovlnné oblasti spekter brzdného záření při napětích 15 kV/1 mA, 25 kV/0,8 mA, 30 kV/0,8 mA, 33 kV/0,8 mA. K měření používejte tyto parametry: clonu o průměru 2 mm, interval Braggova úhlu pro 15 kV v rozmezí (10° – 15°) s krokem 0.2° a dobou expozice 8 s a pro ostatní napětí interval Braggova úhlu (3° – 10°) s krokem 0,2° a dobou expozice 5 s;

        \item proměřte charakteristická spektra rentgenky při napětích 15 kV/1 mA a 33 kV/0,8 mA. K měření používejte tyto parametry: clonu o průměru 2 mm, interval Braggova úhlu (15° – 30°), krok 0,1° a dobu expozice 2 s;

        \item proměřte tvar spektra s Zr absorbérem. K měření používejte tyto parametry: clonu s Zr absorbérem tloušťky 0,05 mm, interval Braggova úhlu (3° – 30°), krok 0,1° a dobu expozice 2 s;

        \item proměřte tvar spektra s Ni absorbérem. K měření používejte tyto parametry: clonu s Ni absorbérem tloušťky 0,01 mm, interval Braggova úhlu (3° – 30°), krok 0,1° a dobu expozice 2 s.
    \end{enumerate}

\item Rentgenka s Fe anodou

    \begin{enumerate}
        \item proměřte charakteristické spektrum rentgenky při napětí 33 kV/0,8 mA. K měření používejte tyto parametry: clonu o průměru 2 mm, interval Braggova úhlu (3° – 30°), krok 0,1° a dobu expozice 2 s;

        \item proměřte tvar spektra s Zr absorbérem. K měření používejte tyto parametry: clonu s Zr absorbérem tloušťky 0,05 mm, interval Braggova úhlu (3° – 30°), krok 0,1° a dobu expozice 3 s.
    \end{enumerate}

\item Rentgenka s Mo anodou.

    \begin{enumerate}
        \item proměřte charakteristické spektrum rentgenky při napětí 33 kV/0.8 mA. K měření používejte tyto parametry: clonu o průměru 2  mm, interval Braggova úhlu (3° – 35°), krok 0,1° a dobu expozice 3 s.
    \end{enumerate}

\item Rentgenka s Cu anodou:

    \begin{enumerate}
        \item proměřte charakteristické spektrum rentgenky při napětí 33 kV/0,8 mA v intervalu Braggova úhlu (42° – 51°). K měření používejte tyto parametry: clonu o průměru 2 mm, krok 0,1° a dobou expozice 2 s.
    \end{enumerate}
    
\end{enumerate}

\item Interpretujte naměřené výsledky (pro mezirovinnou vzdálenost krystalu LiF používejte hodnotu d = 201,4 pm):

\begin{enumerate}
    \item Krátkovlnná mez brzdného záření

    \begin{enumerate}
        \item Ze změřených mezních vlnových délek (respektive frekvencí) určete hodnotu Planckovy konstanty a oceňte přesnost měření
    \end{enumerate}

\item Moseleyův zákon

    \begin{enumerate}
        \item Přesvědčte se, že naměřené úhlové frekvence spektrálních čar K$\alpha$ a K$\beta$ pro různé prvky splňují Moseleyův zákon. Ze směrnice příslušné závislosti určete hodnotu Rydbergovy úhlové frekvence a využitím této hodnoty určete též průměrnou hodnotu stínící konstanty.

        \item Přesvědčte se, že i naměřené polohy absorpčních hran Zr a Ni splňují Moseleyův zákon.

        \item Všimněte si, že absorpční hrana Ni koinciduje se spektrální čarou K$\beta$ mědi; této skutečnosti se využívá v rentgenové difraktografii pro monochromatizaci charakteristického spektra mědi. Z provedeného měření určete filtrační efekt niklu pro čáru K$\beta$.
    \end{enumerate}

\item Úhlová disperze

    \begin{enumerate}
        \item Ze změřených spekter molybdenu určete velikost úhlové disperze pro různé řády difrakce.
    \end{enumerate}
\end{enumerate}

\end{enumerate}

% ----------------------------------------------------------------------
%  Teoretická část
% ----------------------------------------------------------------------
\section{Teoretická část}

Chyba je spočtena pomocí metody přenosu chyb \cite{bib:metoda-prenosu-chyb} jako

\begin{equation}
    \sigma = \sqrt{\sum^n_{i=1} \left( \frac{\partial f}{\partial x_i} \right)^2 \sigma^2_{x_i}}
\end{equation}

% ----------------------------------------------------------------------
%  Výsledky a zpracování měření
% ----------------------------------------------------------------------
\section{Výsledky a zpracování měření}


\subsection{Určení Planckovy konstanty}


\subsection{Moseleyův zákon}

Proměřili jsme charakteristické spektrum pro anody složené z Fe, Mo a Cu při urychlovacím napětím napětí 33 kV a proudu 0,8 mA. Podle zadání $\cite{bib:zadani}$ jsme pro některá měření použili zr absorbér tloušťky 0,05 mm a Ni absorbér s tloušťkou 0,01 mm. Pokud nebyl použit ani jeden z těchto absorbérů měřili jsme se clonou o průměru 2 mm. Doba expozice pro každý úhel byla nastavena na 2 s pro Fe a Cu, pro Mo 3 s. Naměřené výsledky jsou uvedeny v tabulce. Uvedený úhel je úhel detektoru, který je dvojnásobný oproti úhlu dopadajícího svazku, proto tento úhel dále ve výpočtech dělíme 2.



    
% ----------------------------------------------------------------------
%  Diskuse výsledků
% ----------------------------------------------------------------------			
\section{Diskuse výsledků}

% ----------------------------------------------------------------------
%  Závěr
% ----------------------------------------------------------------------
\section{Závěr}
